\documentclass[12pt,letterpaper]{article}
\usepackage[utf8]{inputenc}
\usepackage[spanish,mexico]{babel}
\usepackage[top=2.5cm,bottom=2cm,left=1.8cm,right=2cm]{geometry}
\usepackage{graphicx}
\usepackage{caption}
\usepackage{amssymb}
\usepackage{floatrow}
\usepackage{fancyvrb}
\usepackage{xcolor}
\usepackage{xurl}
\usepackage{pgfgantt}

\PassOptionsToPackage{hyphens}{url}\usepackage{hyperref}
\hypersetup{
	colorlinks=true,
	urlcolor=blue,
	linkcolor=black
}

\pagestyle{headings}

% -- Estructura para insertar imagenes --

%\begin{figure}[H]
%	\centering
%	\includegraphics[width=0.7\linewidth]{img/nombreImg}
%	\caption{Descripción de la imagen.}
%	\label{fig:alias}
%\end{figure}

% --

\begin{document}
	\newgeometry{top=2cm,bottom=2.5cm,left=1cm,right=1cm}
	\begin{titlepage}
		\centering
		\begin{minipage}{0.14\linewidth}
			\includegraphics[width=\linewidth]{img/shieldUnam}
		\end{minipage}
		\begin{minipage}{0.7\linewidth}
			\centering
			{\bfseries\large UNIVERSIDAD NACIONAL AUTÓNOMA DE MÉXICO \par}
			%\vspace{1cm}
			\vfill
			{\scshape\Large Facultad de Ingeniería \par}
			\vfill
			%\vspace{0.7cm}
			{\scshape\Large Ingeniería en Computación \par}
		\end{minipage}
		\begin{minipage}{0.14\linewidth}
			\includegraphics[width=\linewidth]{img/shieldFi}
		\end{minipage}
		
		\centering
		\vspace{1.5cm}
		{\scshape\Large Bases de Datos \par}
		{\scshape\Large grupo: 01\par}
		\vspace{3cm}
		{\scshape\Huge Proyecto final \par}
		\vspace{0.8cm}
		%{\scshape\Huge fff \par}
		\vfill
		{\Large Alumnos: \par}
		\begin{center}
			\begin{tabular}{l}
				$\bullet$ \\
				\\
				$\bullet$ \\
				\\
				$\bullet$ {\Large López González Kevin } \\
				\\
				$\bullet$ \\
				\\
				$\bullet$ \\
			\end{tabular}
		\end{center}
		\vfill
		{\Large Profesor: \par}
		{\Large ING. Fernando Arreola Franco \par}
		\vfill
		{\Large \today \par}
	\end{titlepage}

	\restoregeometry
	\tableofcontents
	\section{Introducción}
	
	\section{Plan de trabajo}
		\subsection{Descripción}
		
		\subsection{Plan de actividades}

		\newpage
		\subsection{Cronograma}
			\begin{ganttchart}[vgrid, x unit=7mm, time slot format=isodate,
				canvas/.style={draw=none},
				title/.append style={fill=blue!10, rounded corners=1mm},
				group/.append style={draw=black, fill=magenta!60},
				bar/.append style={fill=magenta!20}]
				{2021-11-21}{2021-12-09}
				\gantttitlecalendar{year, month=name, day} \\
				
				\ganttgroup{Diseño}{2021-11-21}{2021-11-22}\\
				\ganttbar{Modelo Conceptual}{2021-11-21}{2021-11-21}\\
				\ganttbar{Modelo Lógico}{2021-11-21}{2021-11-21} \\
				\ganttbar{Revisión}{2021-11-22}{2021-11-22} \\
				
				\ganttgroup{Implementación}{2021-11-23}{2021-12-01}\\
				\ganttbar{Modelo Físico}{2021-11-23}{2021-11-25}\\
				\ganttbar{Revisión}{2021-11-26}{2021-11-26}\\
				\ganttbar{Corrección de errores}{2021-11-27}{2021-11-29}\\
				\ganttbar{Pruebas}{2021-11-30}{2021-12-01}\\
				
				\ganttgroup{Presentación}{2021-11-23}{2021-12-01}\\
				\ganttbar{Página Web}{2021-11-23}{2021-11-25}\\
				\ganttbar{Revisión}{2021-11-26}{2021-11-26}\\
				\ganttbar{Corrección de errores}{2021-11-27}{2021-11-29}\\
				\ganttbar{Pruebas}{2021-11-30}{2021-12-01}\\
				
				\ganttgroup{Acoplamiento}{2021-12-02}{2021-12-08}\\
				\ganttbar{Desarrollo}{2021-12-02}{2021-12-04}\\
				\ganttbar{Pruebas}{2021-12-05}{2021-12-05}\\
				\ganttbar{Revisión}{2021-12-06}{2021-12-06}\\
				\ganttbar{Corrección de errores}{2021-12-07}{2021-12-08}\\

				\ganttbar{Documentación}{2021-11-22}{2021-12-08}
			\end{ganttchart}

				
		\subsection{Aportaciones}
		\begin{center}
			\begin{tabular}{c|c|c|c|c|c}
				\hline \\
				& Diseño & Implementación & Presentación & Acoplamiento & Documentación \\
				Kevin López & $\checkmark$ & & $\checkmark$ & $\checkmark$ & $\checkmark$ \\ 
			\end{tabular}
		\end{center}
	
	\section{Diseño}
		\subsection{Análisis de requerimientos}
		
		\subsection{Modelo conceptual}
			\textbf{Entidades}\par 
			\begin{itemize}
				\item PROVEEDOR: \{ \underline{id\_Proveedor}, razón social, domicilio (estado, código postal, colonia, calle y número), nombre, teléfonos \}
				\item CLIENTE: \{\underline{RFC}, nomre (nombre, ap\_Paterno, ap\_Materno), domicilio (estado, código postal, colonia, calle y número), emails \}
				\item INVENTARIO: \{\underline{id\_Inventario}, precio\_compra, fecha\_compra, cantidad\_ejemplares \}
				
				\item PRODUCTO: \{\underline{código\_Barras}, marca, descripción, precio, categoria\}
				\item VENTA : \{\underline{num\_venta}, fecha\_venta, pago\_Total, cantidad\_articulo, pago\_total\_Articulo \}
			\end{itemize}
		
			\textbf{Relaciones}\par
			\begin{itemize}
				\item Un proveedor surte a muchos inventarios.
				\item Un inventario es surtido por muchos proveedores. \\
				\item Un inventario almacena muchos productos.
				\item Un producto es almacenado por un inventario.\\
				\item Una venta contiene muchos productos.
				\item Un producto es contenido es muchas ventas.\\
				\item Un cliente concreta muchas ventas.
				\item Una venta es concretada por un cliente.
				
			\end{itemize}
			\subsubsection{Modelo Entidad-Relación}
			
		\subsection{Modelo lógico}
			\subsubsection{Representación Intermedia}
			\begin{itemize}
				\item PROVEEDOR: \{ id\_proveedor smallint (PK), nombre varchar 50, razón social varchar 50, estado varchar 50, colonia varchar 50, numero smallint, cp smallint, calle varchar 50\}
				\item TELEFONO: \{teléfono bigint(PK), id\_proveedor smallint (FK)\}
				
				\item INVENTARIO: \{id\_Inventario smallint (PK), precio\_compra decimal (10,2), stock smallint, fecha\_compra date \}
				
				\item SURTE: \{[id\_Provedor smaillint (FK), id\_Inventario smallint (FK)] (PK)\}
				
				\item PRODUCTO: \{cod\_barras integer PK, id\_categoria smallint FK, precio smallint NOT NULL, marca varchar(20) NOT NULL, descripcion varchar(50), id\_inventario smallint (FK)\}
				
				\item CATEGORÍA: \{ id\_categoria smallint PK, tipo varchar(20) NOT NULL\}
				
				\item CLIENTE: \{RFC varchar(13) (PK), nombre varchar(20), ap\_paterno varchar (20), ap\_materno varchar (20) (N), cp smallint, numero smallint, estado varchar (32), calle varchar (32), colonia varchar (32)\}
				
				\item EMAIL: \{RFC varchar(13) (FK), email varchar (64) (PK)\}
				
				\item VENTA: \{id\_venta int(PK), fecha\_venta date, pago\_final decimal(7,2), RFC varchar(13)(FK)\}
				
				\item CONTIENE: \{ [cod\_barras int , id\_venta int](PK)(FK), precioTotalArt decimal(7,2), cantidad articulo int\} 
			\end{itemize}
			
			\subsubsection{Transformación de MER a MR}
			
			\subsubsection{Modelo Relacional}
			
			\subsubsection{Normalización}
		
	\section{Implementación}
		\subsection{Modelo físico}
		
			\subsubsection{IaaS}
	
		\subsection{Códigos}
		
		\subsection{DDL}
	
	\section{Presentación}
		\subsection{Página Web}
	
	\section{Conclusiones}
		\begin{itemize}
			\item 
				\subitem 
				
			\item López González Kevin
				\subitem Bla bla bla

			\item 
				\subitem
		\end{itemize}
	
\end{document}



















